\section{Выводы}

Суффиксное дерево эффективно решает задачу множественного поиска за счёт
единовременного построения по тексту. При большом числе запросов это даёт
заметный выигрыш: на случайных данных длиной $30\,000$ и $1000$ образцах
ускорение по сравнению с повторным \texttt{std::string::find} составило
порядка $50\times$, а на тексте длиной $10\,000$ — около $39\times$.
В целом, чем больше запросов, тем выгоднее предварительное построение структуры.

При этом поведение сильно зависит от характера входных данных.
На периодическом тексте и строке из одинаковых символов суффиксное
дерево проигрывает стандартному поиску: при большом числе совпадений
основная стоимость уходит в обход поддерева и обработку результатов,
то делает асимптотику $O(m + k \log k)$ фактически доминирующей по времени.

Отдельно интересен случай редких совпадений. На тесте с уникальными
подстроками (длина $30\,000$) дерево показывает ускорение порядка $38\times$,
так как каждый запрос обрабатывается без повторного прохода по тексту.

Получается, что суффиксное дерево выгодно при большом числе запросов
и умеренном числе совпадений, но теряет эффективность в сценариях с массовыми
вхождениями или когда линейный скан оказывается дешевле константно.

Реализация алгоритма Укконена оказалась сильно сложнее, чем я ожидал:
корректная работа требует аккуратного соблюдения нескольких инвариантов,
и ошибки в них проявляются только на специальных входах с повторениями.
\pagebreak